% ======================================================================
\section{Architecture}
\label{sec:arch}

\begin{figure}[t]
\centering
\begin{tikzpicture}[font=\small, >=Stealth, node distance=4mm,
  box/.style={draw=cMuted, rounded corners=2pt, fill=#1, minimum height=8mm, align=center, inner sep=3pt},
  box/.default=white]
  \node[box=cInkTwo!8] (img) {image\\[-1pt]{\footnotesize $32{\times}32$ / $64{\times}64$}};
  \node[box=cBlue!10, right=6mm of img] (trunk) {frozen CNN trunk\\[-1pt]{\footnotesize trained once on a}\\[-2pt]{\footnotesize disjoint \emph{trunk split}}};
  \node[box=cBlue!4, right=6mm of trunk] (feat) {feature map\\[-1pt]$64\times8\times8$\\[-1pt]{\footnotesize standardised}};
  \draw[->, thick, cInkTwo] (img) -- (trunk);
  \draw[->, thick, cInkTwo] (trunk) -- (feat);
  % tree
  \node[box=cOrange!12, right=14mm of feat, yshift=9mm] (root) {$g_{\mathrm{r}}$};
  \node[box=cOrange!12, below left=8mm and 3mm of root] (A) {$g_{\mathrm{rA}}$};
  \node[box=cOrange!12, below right=8mm and 3mm of root] (B) {$g_{\mathrm{rB}}$};
  \node[box=cOrange!6, below left=6mm and -2mm of A, font=\footnotesize] (AA) {$\cdots$};
  \node[box=cOrange!6, below right=6mm and -2mm of A, font=\footnotesize] (AB) {$\cdots$};
  \node[box=cOrange!6, below right=6mm and -2mm of B, font=\footnotesize] (BB) {$\cdots$};
  \draw[->, thick, cAqua] (root) -- node[left, font=\footnotesize, cInk] {A} (A);
  \draw[->, thick, cOrange] (root) -- node[right, font=\footnotesize, cInk] {B} (B);
  \draw[->, cAqua] (A) -- (AA);
  \draw[->, cOrange] (A) -- (AB);
  \draw[->, cOrange] (B) -- (BB);
  \draw[->, thick, cInkTwo] (feat.east) -- ++(5mm,0) |- (root.west);
  \node[draw=cMuted, dashed, rounded corners, fit=(root)(A)(B)(AA)(AB)(BB), inner sep=4pt,
        label={[font=\footnotesize, cInkTwo]above:tree of small node networks}] {};
  \node[anchor=north west, align=left, font=\footnotesize, text width=8.2cm, cInkTwo] (nodespec)
       at ($(img.south west)+(0,-7mm)$)
  {node $g_v$: $1{\times}1$ conv ($64\to k$) $\to$ ReLU $\to$ linear ($64k\to w$) $\to$ ReLU $\to$ linear ($w\to1$),
   $k=\lceil w/4\rceil$; $w=8$ gives 1{,}171 parameters};
\end{tikzpicture}
\caption{The model. A small CNN trunk, trained once on data the tree never sees, maps every image to a
$64\times8\times8$ feature map. Every node of the tree is a small neural network on that map. Child A
corrects the parent's misses and child B its false alarms.}
\label{fig:arch}
\end{figure}

\subsection{Features: one shared trunk, trained separately}
\label{sec:trunk}

Running a convolutional network inside every node would multiply the cost of the convolutions by the
number of nodes. Instead, every node reads the same frozen feature map produced by a small CNN, the
\emph{trunk}. This is the shared-trunk design proposed in Section~4.1 of the first report~\cite{saini2026tree}. The trunk
(\cref{fig:arch}) has two VGG-style stages (two $3\times3$ conv--BN--ReLU layers each, then max-pooling; a
stride-2 stem is added for the $64\times64$ Cats vs Dogs images), giving a $64\times8\times8$ map. For
pre-training only, a third stage and a linear classifier are added on top; the nodes never see them. The tap
therefore exposes \emph{mid-level} features, so the tree still has real work to do.

The key design decision is \emph{which data the trunk is trained on}. Each dataset is split into four
disjoint parts (\cref{tab:splits}). The trunk is trained only on the \emph{trunk split}, and the tree is
grown on the separate \emph{tree split}. If the trunk had seen the tree's training images, their features
would be unnaturally easy: the trunk has already memorised them (its own training accuracy is 97--99\%).
The errors that drive the tree's growth would then be unrepresentative of the errors made on new images.
This is the same reason stacked ensembles use out-of-fold predictions~\cite{wolpert1992stacked}. Three trunks are compared:
\begin{itemize}
\item \textbf{own split}: supervised on the dataset's trunk split (40 epochs, SGD, crop/flip augmentation);
\item \textbf{random}: the same architecture left untrained, with only its BatchNorm statistics calibrated, so
the tree must do all the learning;
\item \textbf{transfer} (Cats vs Dogs only): the CIFAR-10 trunk applied to images resized to $32\times32$. CIFAR-10
contains cats and dogs, but at a different resolution and from different sources.
\end{itemize}
Features are standardised per channel and cached in half precision, so a node is trained entirely on cached
tensors.

\begin{table}[t]
\centering
\small
\caption{Disjoint data splits. Cats vs Dogs is Microsoft's Kaggle release of the Asirra photos~\cite{elson2007asirra} (25{,}000 photos;
three undecodable files dropped); CIFAR-10 is from~\cite{krizhevsky2009learning}, centre-cropped and resized to $64\times64$.}
\label{tab:splits}
\begin{tabular}{@{}l r r r r r@{}}
\toprule
dataset & trunk & tree (train) & val & test & classes\\
\midrule
CIFAR-10 & 15{,}000 & 30{,}000 & 5{,}000 & 10{,}000 & 10\\
Cats vs Dogs & 6{,}000 & 11{,}997 & 2{,}000 & 5{,}000 & 2\\
\bottomrule
\end{tabular}
\end{table}

\subsection{Nodes: small neural networks}
\label{sec:node}

A node of width $w$ is the network
\[
g_v(f)=W_3\,\mathrm{ReLU}\big(W_2\,\mathrm{vec}(\mathrm{ReLU}(W_1 * f))\big),\qquad
W_1:\ 1{\times}1\text{ conv } 64\to k,\quad W_2: 64k\to w,\quad W_3: w\to1,
\]
with $k=\lceil w/4\rceil$ and biases throughout. The $1\times1$ convolution chooses which trunk channels matter,
and the linear layer over all 64 spatial positions decides \emph{where} in the image they matter. That
spatial weighting is what the receptive-field analysis of \cref{sec:rf} visualises. At the default width $w=8$
a node has 1{,}171 parameters. Nodes are trained with Adam on a class-balanced logistic loss, which gives a
good ranking even when a deep child has only a handful of positives among thousands of samples. Two further
steps follow the original LP formulation, which minimises the number of misclassified samples. After every
pass over the data, the output threshold is moved to the value that minimises the node's training errors,
never letting it become a constant. Training stops when the node makes no errors on its own sample set, or
after 20 passes without improvement.

\subsection{Growing the tree}
\label{sec:grow}

Node $v$ owns samples $S_v$ with binary targets $t_v$. After fitting $g_v$, its samples fall into the four
categories of the original paper~\cite{jayadeva2021tree}: C1 ($g=0,t=0$), C2 ($g=1,t=1$), C3 ($g=0,t=1$, misses) and C4 ($g=1,t=0$, false
alarms). Child A must turn the misses into 1s and child B the false alarms into 0s. We implement both
constructions of the children.

\paragraph{Faithful (as in the first assignment).} A is trained on \emph{all} of $S_v$, with target 1 on C3 and
0 elsewhere. B is trained on all of $S_v$, with target 1 on C4 and 0 elsewhere. The output is
$y=(g\lor A)\land\lnot B$.

\paragraph{Gated (exploiting the don't-cares).} Table~II of the paper marks A's output on C2 and B's output on
C1 as ``don't care''. In fact, under the override rule, A's output matters \emph{only} where $g=0$ and B's only
where $g=1$. We therefore train A on $\{g=0\}=C1\cup C3$ with the true targets, and B on $\{g=1\}=C2\cup C4$ with
flipped targets, and predict
\[
y(x)=\begin{cases}A(x) & \text{if } g(x)=0,\\ \lnot B(x) & \text{if } g(x)=1.\end{cases}
\]
This changes the algorithm in three useful ways.
\begin{enumerate}[(i)]
\item The children's sample sets \emph{partition} $S_v$. Each level of the tree costs one pass over the data in
total, instead of one pass per node, which was Challenge~2 of the first report.
\item A child separates its targets only from the samples the parent sent its way, not from the whole dataset.
This is an easier problem with milder class imbalance (Challenge~3).
\item Termination is immediate. A node that is not constant sends strictly fewer samples to each child, so
every path ends after at most $|S_v|$ levels, and the leaves classify their samples perfectly unless two
identical inputs carry different labels.
\end{enumerate}
Both constructions agree on how a fully grown tree classifies its training data. They differ in which subproblems
the children are asked to solve, which turns out to matter a great deal (\cref{sec:results}). A node whose sample set is pure gets a constant unit with no network. Trees are grown until every node is
error-free or a depth cap of 40 is reached.

\subsection{Multiclass: a tree of class groups}
\label{sec:hier}

For CIFAR-10 we implement the hierarchical class-grouping scheme proposed in Section~5.3 of the first
report~\cite{saini2026tree},
rather than ten one-vs-rest trees. The trunk's confusion matrix on the validation split defines a graph whose
edge weights are how often two classes are confused. We split the classes recursively along the Fiedler vector
of its normalised Laplacian~\cite{shi2000normalized,vonluxburg2007tutorial}, choosing the cut with the smallest normalised cut and requiring at least two
classes per side when possible. The 9 internal splits of the resulting binary class hierarchy are each solved
by a binary tree-type network, trained only on the samples of the classes below that split. An image is
classified by routing it down the hierarchy. Because every training image is routed correctly at every level
once every split reaches 100\% training accuracy, the whole multiclass tree then reaches 100\% training
accuracy. Classes are never compared against an unrelated jumble of other classes, and each split's tree only
has to solve the confusion that actually occurs there.
